\chapter{Optimization}\label{ch:optimize}

\laheader{optimize.h, optimize.cc}

The optimization module provides local minimization of a function $f:\mathbb{R}^n\to\mathbb{R}$ (and, for two methods, the
search for saddle points) by a family of methods sharing one interface: quasi-Newton methods with line search (BFGS, DFP, SR1,
L-BFGS), nonlinear conjugate gradients, GDIIS, trust-region Newton and quasi-Newton methods, partitioned rational function
optimization (P-RFO) for saddle points, Levenberg-Marquardt for least squares, the derivative-free methods of Nelder and Mead,
Powell and COBYLA (with nonlinear inequality constraints), and the first-order and stochastic methods SGD, Adam, SPSA, the
quantum natural gradient and Rotosolve used in machine learning and for variational quantum algorithms.

The design separates three things:
\begin{itemize}
\item the \emph{problem}: a class derived from \code{OptProblem<T>} implementing \code{funcgrad()}, optionally also
deriving from \emph{capability} interfaces (Hessian, metric, trigonometric structure, constraints, residuals) which some
methods require or exploit (section~\ref{sec:optimize:problem});
\item the \emph{method}: a class derived from \code{Optimizer<T>} with the single virtual function \code{minimize()};
\item the \emph{hyperparameters}: plain value objects of a class hierarchy rooted in \code{hyperparameters<T>}, one derived
class per method, passed to \code{minimize()} and returned filled with the values actually reached
(section~\ref{sec:optimize:hyper}).
\end{itemize}
Methods are written on top of a small infrastructure (evaluation counter, iteration bookkeeping, line searches) and of three
common drivers (line search, first order, trust region), so that a new method usually consists of a short class supplying the
search direction, the update rule or the trust-region step (section~\ref{sec:optimize:drivers}).

All templates are parameterized by the real type \code{T}; optimize.cc instantiates them for \code{T = double} and
\code{float}. The capability interfaces, adapters and the hyperparameter classes are header-only templates; the optimizers, the
adapters with member functions in optimize.cc (\code{FrozenOptProblem}, \code{ParameterShift}, \code{FiniteDifference},
\code{OptLeastSquares}) and the auxiliary functions exist in the library for \code{double} and \code{float}.
The common drivers \code{linesearch_minimize}, \code{firstorder_minimize} and \code{trustregion_minimize}, the extension
points for user-written methods, are explicitly instantiated for both types as well. The defaults depending on the precision
are taken from \code{opt_precision<T>} (section~\ref{sec:optimize:hyper}).

The data members of all classes of this chapter are small, except vectors held by \code{FrozenOptProblem} and
\code{optevaluator}; these are ordinary \code{NRVec} objects following the reference counting and copy-on-demand scheme of
section~\ref{sec:concepts:refcount}, which has consequences for user code described in section~\ref{sec:optimize:pitfalls}.
The library code of this module calls \code{copyonwrite()} (or works on freshly allocated vectors) before every element write;
the only weak point found is noted in section~\ref{sec:optimize:adapters}.

%-----------------------------------------------------------------------------------------------------------------------------
\section{Problem interface and capabilities}\label{sec:optimize:problem}

\subsection{OptProblem}

\begin{lareference}
\lafunc{template<typename T> class OptProblem}{Abstract base class of every optimization problem. A user problem derives from it
(and possibly from capability classes) and stores any auxiliary data needed to evaluate $f$.}
\lafunc{virtual T funcgrad(const NRVec<T> &x, NRVec<T> *grad) const =0}{Returns $f(x)$. If \code{grad!=NULL}, also stores
$\nabla f(x)$ in \code{*grad}. The gradient vector is allocated by the caller with \code{x.size()} elements, but it may share
its data with another vector, so the implementation must call \code{grad->copyonwrite()} before writing its elements (or
assign a whole new vector to \code{*grad}). Derivative-free methods always pass \code{NULL}; a problem not able to provide the
gradient may then simply ignore \code{grad} (or call \code{laerror} if it is not \code{NULL}) and be wrapped in
\code{FiniteDifference} or \code{ParameterShift} for gradient methods.}
\end{lareference}

\subsection{Capabilities}\label{sec:optimize:capabilities}

Further information about the problem is supplied by deriving the problem class also from one or more of the following
abstract classes (multiple inheritance). Methods look them up at run time by \code{optcapability<C>(problem)}; a method
requiring a capability calls \code{laerror} at the start if it is missing, a method that can only use it optionally falls back
silently (with a verbose message in some cases). Matrices and vectors passed to the capability functions are allocated by the
caller with the right dimensions; implementations writing their elements have to call \code{copyonwrite()} first (or assign
whole objects, e.g.\ \code{h=(T)0;} followed by element writes, since assignment of a scalar is a modifying member function).

\begin{lareference}
\lafunc{template<typename T> class OptHessian}{Hessian of $f$. Used by \code{TrustRegion_Optimizer} and
\code{PRFO_Optimizer}.}
\lafunc{virtual void hessian(const NRVec<T> &x, NRSMat<T> &h) const =0}{Hessian at \code{x} into \code{h} ($n\times n$
symmetric, allocated by the caller).}
\lafunc{virtual NRVec<T> hessian_times(const NRVec<T> &x, const NRVec<T> &v) const}{Hessian-vector product $H(x)\,v$; the
default builds the full Hessian by \code{hessian()} and multiplies. It can be overridden for large problems.
Currently no method calls it (the Steihaug subproblem works with the stored matrix $B$; \code{FrozenOptProblem} forwards it);
it is kept for future Hessian-free methods.}

\lafunc{template<typename T> class OptMetric}{Metric tensor of the parameter space: the Fubini-Study metric for the quantum
natural gradient \cite{optimize:stokes2020} or the Fisher information for the natural gradient in machine
learning \cite{optimize:amari1998}. Used by \code{QNG_Optimizer} and by every first-order method with \code{my_natural=true}.}
\lafunc{virtual void metric(const NRVec<T> &x, NRSMat<T> &g) const =0}{Metric $G(x)$ into \code{g} ($n\times n$).}

\lafunc{template<typename T> class OptTrigonometric}{Declares that $f$ as a function of the single parameter $x_i$ is a
trigonometric polynomial $a_0+\sum_{k=1}^{R_i}\bigl(a_k\cos(k\omega_i x_i)+b_k\sin(k\omega_i x_i)\bigr)$,
$\omega_i=2\pi/\mathrm{period}_i$. For parameterized quantum circuits with gates $\exp(-\mathrm{i}x_iP/2)$, $P$ a Pauli string,
$R_i=1$ and the period is $2\pi$. Used by \code{Rotosolve_Optimizer} and \code{ParameterShift}.}
\lafunc{virtual int degree(const int i) const =0}{Degree $R_i$; 0 means that $f$ is not trigonometric in $x_i$ (such
parameters are treated by line minimization or central differences).}
\lafunc{virtual T period(const int i) const}{Period of $f$ in $x_i$; default $2\pi$.}

\lafunc{template<typename T> class OptConstraints}{Inequality constraints $c_i(x)\ge 0$, $i=0,\dots,m-1$, used by
\code{Cobyla_Optimizer}. An equality constraint is given as two inequalities $c(x)\ge0$, $-c(x)\ge0$. Other methods ignore
constraints.}
\lafunc{virtual int nconstraints() const =0}{Number $m$ of constraints.}
\lafunc{virtual void constraints(const NRVec<T> &x, NRVec<T> &c) const =0}{Values $c_i(x)$ into \code{c}, allocated by the
caller with \code{nconstraints()} elements.}

\lafunc{template<typename T> class OptResiduals}{Least-squares problem $f(x)=\frac12|r(x)|^2$ given by the residual vector
$r\in\mathbb{R}^m$ and optionally its Jacobian. Required by \code{LM_Optimizer}.}
\lafunc{virtual int nresiduals() const =0}{Number $m$ of residuals.}
\lafunc{virtual void residuals(const NRVec<T> &x, NRVec<T> &r) const =0}{Residuals into \code{r} (allocated by the caller with
\code{nresiduals()} elements).}
\lafunc{virtual bool jacobian(const NRVec<T> &x, NRMat<T> &J) const}{Jacobian $J_{ij}=\partial r_i/\partial x_j$ into
\code{J} ($m\times n$, allocated by the caller); returns \code{true} if provided. The default returns \code{false}, and the
callers then use forward differences.}
\end{lareference}

\begin{lareference}
\lafunc{template<class C, typename T> const C *optcapability(const OptProblem<T> *o)}{Returns the capability \code{C}
(e.g.\ \code{OptHessian<T>}) of the problem, or \code{NULL}. It is found by \code{dynamic_cast} on the problem itself and, if
the problem is an \code{OptProblemAdapter}, recursively on the wrapped problem. For an adapter with
\code{transforms_capabilities()==true} the adapter's own implementation of \code{C} is returned only if the wrapped problem also
provides \code{C}, and the untransformed capability of the wrapped problem is never passed on. Methods always use this function
(not a plain \code{dynamic_cast}), so that capabilities remain available through adapters.}
\end{lareference}

\subsection{Least-squares problems: OptLeastSquares}

\begin{lareference}
\lafunc{template<typename T> class OptLeastSquares : public OptProblem<T>, public OptResiduals<T>}{Convenience base class for
least-squares problems. The user derives from it and implements \code{nresiduals()}, \code{residuals()} and optionally
\code{jacobian()}. Its \code{funcgrad()} provides $f=\frac12|r|^2$ and $\nabla f=J\Tr r$, so that every method, not only
\code{LM_Optimizer}, can be applied to the problem.}
\lafunc{T fd_step}{Relative step $h$ of the forward-difference Jacobian used when \code{jacobian()} returns \code{false}.}
\lafunc{OptLeastSquares(const T h=1e-7)}{Sets \code{fd_step}.}
\lafunc{void jacobian_fd(const NRVec<T> &x, const NRVec<T> &r, NRMat<T> &J) const}{Jacobian by forward differences, columns
$(r(x+h_je_j)-r(x))/h_j$ with $h_j=h\max(1,|x_j|)$; \code{r} are the residuals at \code{x} ($n$ further residual evaluations).}
\lafunc{virtual T funcgrad(const NRVec<T> &x, NRVec<T> *grad) const}{$\frac12|r(x)|^2$; with \code{grad}, the gradient $J\Tr r$
(from \code{jacobian()} or \code{jacobian_fd()}).}
\end{lareference}
Note that \code{LM_Optimizer} does not call \code{funcgrad()}: it evaluates \code{residuals()} directly (counted as function
evaluations) and uses \code{LM_hyperparameters::my_fd_step}, not \code{fd_step}, for its own finite-difference Jacobian.

\subsection{Adapters}\label{sec:optimize:adapters}

An adapter is a problem wrapping another one, e.g.\ to supply gradients computed from function values or to optimize in a
subspace. Adapters can be nested.

\begin{lareference}
\lafunc{template<typename T> class OptProblemAdapter : public OptProblem<T>}{Base class of adapters; \code{funcgrad()} is left to
the derived class.}
\lafunc{const OptProblem<T> *p}{The wrapped problem (not owned).}
\lafunc{OptProblemAdapter(const OptProblem<T> *pp)}{Wraps \code{pp}.}
\lafunc{virtual bool transforms_capabilities() const}{\code{false} (default): the capabilities of the wrapped problem are
accessible unchanged through \code{optcapability()}. \code{true}: the adapter provides them in a transformed form (e.g.\ in a
different dimension) by its own implementations, see \code{optcapability()}.}
\end{lareference}

\paragraph{FiniteDifference.}
\begin{lareference}
\lafunc{template<typename T> class FiniteDifference : public OptProblemAdapter<T>}{Gradient by central differences
$(f(x+he_i)-f(x-he_i))/(2h)$, i.e.\ $2n$ evaluations of the wrapped problem besides $f(x)$. Makes all gradient methods usable
for a problem providing only $f$. Capabilities of the wrapped problem pass through.}
\lafunc{T h}{Difference step (absolute).}
\lafunc{FiniteDifference(const OptProblem<T> *pp, const T hh=1e-6)}{Wraps \code{pp}.}
\lafunc{virtual T funcgrad(const NRVec<T> &x, NRVec<T> *grad) const}{$f$ of the wrapped problem; the gradient by central
differences if requested.}
\end{lareference}

\paragraph{ParameterShift.}
\begin{lareference}
\lafunc{template<typename T> class ParameterShift : public OptProblemAdapter<T>}{Exact gradient of a problem providing
\code{OptTrigonometric} by the general parameter-shift rule \cite{optimize:wierichs2022} (for degree 1 the classical two-point
rule \cite{optimize:mitarai2018,optimize:schuld2019}): for $x_i$ of degree $R$,
$\partial_i f=\omega\sum_{\mu=1}^{2R}(-1)^{\mu-1}f(x+s_\mu e_i/\omega)/(4R\sin^2(s_\mu/2))$, $s_\mu=(2\mu-1)\pi/(2R)$, i.e.\
$2R_i$ evaluations per parameter. Parameters of degree 0 are differentiated by central differences with step \code{h}.
Capabilities of the wrapped problem pass through.}
\lafunc{const OptTrigonometric<T> *trig}{The trigonometric capability of the wrapped problem.}
\lafunc{T h}{Central-difference step for parameters of degree 0.}
\lafunc{ParameterShift(const OptProblem<T> *pp, const T hh=1e-6)}{Wraps \code{pp}; \code{laerror} if it does not provide
\code{OptTrigonometric}.}
\lafunc{virtual T funcgrad(const NRVec<T> &x, NRVec<T> *grad) const}{$f$ of the wrapped problem and the parameter-shift
gradient.}
\end{lareference}

\paragraph{FrozenOptProblem.}
Optimization in a subset of the parameters, the others (``frozen'') keeping their values.
\begin{lareference}
\lafunc{template<typename T> class FrozenOptProblem : public OptProblemAdapter<T>, public OptHessian<T>, public OptMetric<T>,
public OptTrigonometric<T>, public OptConstraints<T>}{Problem in the $n_{\rm free}$ free parameters. Its
\code{transforms_capabilities()} returns \code{true}: the Hessian and metric (submatrices for the free parameters), the
Hessian-vector product, degree and period (of the free parameters) and the constraints are provided if and only if the wrapped
problem provides them. \code{OptResiduals} is not forwarded, so \code{LM_Optimizer} cannot be used through it (the other
methods can, via \code{funcgrad()} of an \code{OptLeastSquares} problem).}
\lafunc{NRVec<T> xfull}{Full parameter vector; the values of the frozen parameters are taken from it.}
\lafunc{NRVec<int> freeindex}{Indices of the free parameters in the full vector (ascending).}
\lafunc{FrozenOptProblem(const OptProblem<T> *pp, const NRVec<T> &x, const bitvector &frozen)}{Wraps \code{pp}; \code{xfull}
shares its data with \code{x} (shallow copy); \code{frozen} has bit $i$ set if $x_i$ is frozen.}
\lafunc{void setfrozen(const bitvector &frozen)}{New set of frozen parameters (\code{laerror} if
\code{frozen.size()!=xfull.size()}).}
\lafunc{void setfull(const NRVec<T> &x)}{New full vector (new values of the frozen parameters); shallow copy.}
\lafunc{NRVec<T> reduce(const NRVec<T> &full) const}{The free parameters of a full vector, \code{full.subvector(freeindex)}.}
\lafunc{NRVec<T> expand(const NRVec<T> &reduced) const}{Full vector: a copy of \code{xfull} with the free parameters replaced by
\code{reduced}.}
\lafunc{virtual T funcgrad(const NRVec<T> &x, NRVec<T> *grad) const}{$f$ at \code{expand(x)}; the gradient is the free part of
the full gradient.}
\lafunc{virtual void hessian(const NRVec<T> &x, NRSMat<T> &h) const; virtual NRVec<T> hessian_times(const NRVec<T> &x, const NRVec<T> &v) const}{Free-free block of the wrapped Hessian; product with the free block.}
\lafunc{virtual void metric(const NRVec<T> &x, NRSMat<T> &g) const}{Free-free block of the wrapped metric.}
\lafunc{virtual int degree(const int i) const; virtual T period(const int i) const}{Of the free parameter $i$, i.e.\ of
\code{freeindex[i]} of the wrapped problem.}
\lafunc{virtual int nconstraints() const; virtual void constraints(const NRVec<T> &x, NRVec<T> &c) const}{Constraints of the
wrapped problem at \code{expand(x)}.}
\end{lareference}
Usage: \code{FrozenOptProblem<T> fp(&problem,x,frozen); NRVec<T> xr=fp.reduce(x); optimizer.minimize(&fp,xr); x=fp.expand(xr);}
(see example~\ref{sec:optimize:ex:frozen}).

%-----------------------------------------------------------------------------------------------------------------------------
\section{Optimizers, hyperparameters and results}\label{sec:optimize:hyper}

\subsection{Optimizer and optresult}

\begin{lareference}
\lafunc{template<typename T> class Optimizer}{Abstract base class of all methods; the methods can be used through
\code{Optimizer<T> *} or \code{&}. Each method class (\code{BFGS_Optimizer} etc.) has no data members and a virtual destructor.}
\lafunc{virtual optresult<T> minimize(const OptProblem<T> *o, NRVec<T> &x, const hyperparameters<T> *par=NULL, hyperparameters<T> *achieved=NULL) =0}{Minimizes
\code{o} (for P-RFO: searches a stationary point of the requested order) starting from \code{x}, which is replaced by the
result. \code{par==NULL} means the default parameters of the method; otherwise the parameters are taken from \code{*par} as
described in section~\ref{sec:optimize:setfrom}. If \code{achieved!=NULL}, it is filled with the parameters used and with the
values actually reached (section~\ref{sec:optimize:achieved}).}
\lafunc{template<typename T> class optresult}{Result of \code{minimize()}.}
\lafunc{T f; bool converged}{Final function value (for SPSA an estimate); whether a convergence criterion was satisfied
(\code{false} if the iterations stopped on \code{my_maxstep}, \code{my_maxfunc} or a failure).}
\lafunc{optresult(const T ff, const bool c); operator T() const}{Constructor; implicit conversion to \code{f}, so
\code{double f=o.minimize(&p,x);} works.}
\end{lareference}

\subsection{Hyperparameter hierarchy}

The hyperparameters form a class hierarchy; each method has its own class adding its specific parameters to those of the
general classes. All members are public, named \code{my_...}, and initialized to the method's defaults by the constructor. The
hierarchy (defaults in parentheses; a derived class may change inherited defaults, given at the method):

\begin{lstlisting}
hyperparameters                    maxstep(200) maxfunc(0) f_thr(0) x_thr(0) saddle_order(0) verbose(1)
 +- gradient_hyperparameters       g_thr(1e-6)
 |   +- linesearch_hyperparameters    c1(1e-4) c2(0.9) maxlinesearch(30) maxsteplength(1e10)
 |   |   +- quasinewton_hyperparameters   sy_thr(1e-12)
 |   |   |   +- BFGS_hyperparameters         update(QN_BFGS) sr1_thr(1e-8)
 |   |   |   |   +- GDIIS_hyperparameters       diis_dim(5) diis_eps(0)
 |   |   |   +- LBFGS_hyperparameters        memory(10)
 |   |   +- CG_hyperparameters            beta(CG_PR) restart(0) powell_thr(0.2); c2=0.1
 |   +- firstorder_hyperparameters    lr(1e-3) lr_decay(0) lr_power(1) weight_decay(0) return_best(false)
 |   |   |                            natural(false) metric_reg(1e-3); maxstep=1000
 |   |   +- SGD_hyperparameters           momentum(0.9) nesterov(false); lr=1e-2
 |   |   |   +- QNG_hyperparameters          momentum=0 natural=true lr=0.05
 |   |   +- Adam_hyperparameters          beta1(0.9) beta2(0.999) eps(1e-8) amsgrad(false)
 |   |   +- SPSA_hyperparameters          c(0.1) gamma(0.101) resamplings(1) calibrate(20) target_step(0.1);
 |   |                                    lr_decay=0.01 lr_power=0.602 g_thr=0
 |   +- trustregion_hyperparameters   delta0(1) deltamax(1e3) eta(1e-4) hessian(TR_AUTO) hessian_recompute(1)
 |   |   |                            update_thr(1e-8)
 |   |   +- TrustRegion_hyperparameters   subproblem(TR_EIGEN) maxcg(0)
 |   |   +- PRFO_hyperparameters          delta0=0.3
 |   +- LM_hyperparameters            tau(1e-3) scaling(false) fd_step(1e-7)
 +- NelderMead_hyperparameters     initstep(0.1) adaptive(true) alpha(1) gamma(2) rho(0.5) sigma(0.5);
 |                                 maxstep=5000 f_thr=1e-8 x_thr=1e-6
 +- Powell_hyperparameters         initstep(0.1) linmin_tol(1e-4) maxlinemin(50); maxstep=1000 f_thr=1e-10 x_thr=1e-6
 |   +- Rotosolve_hyperparameters     reevaluate(false)
 +- Cobyla_hyperparameters         rhobeg(0.5) rhoend(1e-6) flexible_radius(true); maxstep=100000
\end{lstlisting}

The defaults given above and below are those for \code{double}. The precision-dependent ones are taken from the trait
\code{opt_precision<T>} (\header{optimize.h}, analogous to \code{LA_precision<T>} of section~\ref{sec:concepts:traits}), defined
for \code{T = double} and \code{float} only; the double values are those used by the library before. In float the gradient norm
reachable by a line search is limited by the resolution of $f$ (about $\epsilon|f|$ divided by the step), hence the larger
\code{g_thr}.
\begin{center}
\begin{tabular}{llll}
\toprule
member & double & float & default of \\
\midrule
\code{fd_central_step()} & $10^{-6}$ & $5\cdot10^{-3}$ & \code{hh} of \code{FiniteDifference}, \code{ParameterShift} \\
\code{fd_forward_step()} & $10^{-7}$ & $3\cdot10^{-4}$ & \code{h} of \code{OptLeastSquares}, \code{my_fd_step} (LM) \\
\code{g_thr()} & $10^{-6}$ & $10^{-3}$ & \code{my_g_thr} \\
\code{sy_thr()} & $10^{-12}$ & $10^{-6}$ & \code{my_sy_thr} (quasi-Newton) \\
\code{sr1_thr()} & $10^{-8}$ & $10^{-4}$ & \code{my_sr1_thr} (SR1) \\
\code{update_thr()} & $10^{-8}$ & $10^{-4}$ & \code{my_update_thr} (trust region) \\
\code{nm_f_thr()}, \code{nm_x_thr()} & $10^{-8}$, $10^{-6}$ & $10^{-5}$, $10^{-4}$ & \code{my_f_thr}, \code{my_x_thr} (Nelder-Mead) \\
\code{powell_f_thr()}, \code{powell_x_thr()} & $10^{-10}$, $10^{-6}$ & $10^{-6}$, $10^{-4}$ & \code{my_f_thr}, \code{my_x_thr} (Powell) \\
\code{rhoend()} & $10^{-6}$ & $10^{-4}$ & \code{my_rhoend} (COBYLA) \\
\bottomrule
\end{tabular}
\end{center}

\paragraph{Base classes.}
\begin{lareference}
\lafunc{template<typename T> class hyperparameters}{Parameters common to all methods.}
\lafunc{int my_maxstep}{Maximum number of iterations (default 200). What an iteration is depends on the method (a line search
step, a sweep over all directions or parameters, a simplex operation, a trust-region step including rejected ones).}
\lafunc{int my_maxfunc}{Maximum number of function evaluations, $\le0$ = unlimited (default 0). It counts every call of
\code{funcgrad()} by the method (with or without gradient), and for LM every residual evaluation, including those of
finite-difference Jacobians. Evaluations inside adapters (finite differences, parameter shifts) are not counted. Methods check
before each iteration that the budget suffices and never exceed it, except for the unavoidable initial evaluations (the
starting point, the initial simplex of Nelder-Mead and COBYLA, the first SPSA estimate).}
\lafunc{T my_f_thr; T my_x_thr}{Thresholds on $|\Delta f|$ and $|\Delta x|$ (Euclidean norm) in the last step (default 0).
The general criterion is $|\Delta f|<$ \code{my_f_thr} \emph{and} $|\Delta x|<$ \code{my_x_thr}, both strict, so with the
default 0 it is switched off, and setting only one of them has no effect. Some methods reinterpret them (Nelder-Mead, Powell,
Rotosolve, COBYLA, see there).}
\lafunc{int my_saddle_order}{0 = minimum (default), $k>0$ = saddle point with $k$ negative Hessian eigenvalues. Only
\code{PRFO_Optimizer} accepts $k>0$; all other methods call \code{laerror} if it is nonzero. For the trust-region methods it
returns the Hessian index in \code{achieved}.}
\lafunc{int my_verbose}{0 = silent, 1 = one line per iteration and a final report (default), 2 = more details (line searches,
step lengths, trust-region radii, restarts). Output goes to \code{std::cout}.}
\lafunc{hyperparameters()}{Defaults as above.}
\lafunc{virtual void setfrom(const hyperparameters &p)}{Copies the parameters from \code{p} which exist in both classes
(section~\ref{sec:optimize:setfrom}).}
\lafunc{virtual bool converged(const optstate<T> &s) const}{Convergence test: \code{s.stepdone && s.df<my_f_thr && s.dx<my_x_thr}.}
\lafunc{virtual void setreached(const optstate<T> &s)}{Stores the reached values: \code{my_maxstep} = iterations done,
\code{my_maxfunc} = evaluations done, \code{my_f_thr}, \code{my_x_thr} = $|\Delta f|$, $|\Delta x|$ of the last step.}

\lafunc{template<typename T> class gradient_hyperparameters : public hyperparameters<T>}{Methods using the gradient.}
\lafunc{T my_g_thr}{Convergence threshold on the gradient norm (default \code{1e-6}). \code{converged()} is
$|g|<$ \code{my_g_thr} \emph{or} the criterion of the base class; \code{setreached()} stores the last $|g|$ (or $-1$ if the
gradient norm is not available, as for SPSA).}

\lafunc{template<typename T> class linesearch_hyperparameters : public gradient_hyperparameters<T>}{Methods with the strong
Wolfe line search.}
\lafunc{T my_c1; T my_c2}{Sufficient decrease (Armijo) and curvature parameters of the strong Wolfe conditions,
$0<c_1<c_2<1$ (defaults \code{1e-4}, 0.9).}
\lafunc{int my_maxlinesearch}{Maximum number of function evaluations in one line search (30).}
\lafunc{T my_maxsteplength}{Maximum norm of a single step in $x$ (\code{1e10}).}

\lafunc{template<typename T> class firstorder_hyperparameters : public gradient_hyperparameters<T>}{First-order methods with a
learning rate instead of a line search, suitable also for noisy functions and gradients. Default \code{my_maxstep=1000}.}
\lafunc{T my_lr; T my_lr_decay; T my_lr_power}{Learning rate in iteration $k=0,1,\dots$ is
$\mathrm{lr}/(1+\mathrm{lr\_decay}\cdot k)^{\mathrm{lr\_power}}$ (defaults \code{1e-3}, 0, 1).}
\lafunc{T my_weight_decay}{Decoupled weight decay $x\mathrel{-}=\mathrm{lr}\cdot\mathrm{weight\_decay}\cdot x$
\cite{optimize:loshchilov2019} (default 0).}
\lafunc{bool my_return_best}{Return the point with the lowest function value seen instead of the last one (default
\code{false}); only meaningful for (nearly) deterministic $f$, the lowest of noisy values is biased.}
\lafunc{bool my_natural}{Use the natural gradient $(G+\mathrm{metric\_reg}\,I)^{-1}g$ with the metric $G$ of the problem
(\code{OptMetric} required) \cite{optimize:amari1998,optimize:stokes2020}; $|g|$ in the convergence test and in the output then
refers to the natural gradient (default \code{false}).}
\lafunc{T my_metric_reg}{Regularization added to the diagonal of the metric (\code{1e-3}).}

\lafunc{template<typename T> class trustregion_hyperparameters : public gradient_hyperparameters<T>}{Trust-region methods,
see section~\ref{sec:optimize:tr}.}
\end{lareference}

The method-specific classes are described with their methods in section~\ref{sec:optimize:methods}. Each of them has a default
constructor and overrides \code{setfrom()}.

\subsection{How parameters are taken: setfrom()}\label{sec:optimize:setfrom}

Every \code{minimize()} constructs the hyperparameter object of its own method (with the method's defaults) and calls
\code{hp.setfrom(*par)}. Each class's \code{setfrom()} tries \code{dynamic_cast} of \code{par} to its own type; on success it
assigns \code{*this=*q} at its own level, otherwise it delegates to its parent class. The effect is that all fields of the
deepest class common to the method's class and the dynamic type of \code{*par} are copied from \code{*par}, while the
method-specific fields below it keep the method's defaults. Thus one may pass
\begin{itemize}
\item the method's own class (all parameters given),
\item a more general class, e.g.\ \code{gradient_hyperparameters<double>} to any gradient method, or \code{hyperparameters<double>}
to any method,
\item the class of another method: e.g.\ \code{CG_hyperparameters} passed to \code{BFGS_Optimizer} provides the line search
parameters (including CG's default \code{my_c2=0.1}).
\end{itemize}
\lawarning{All fields of the common level are copied, including those the caller left at the defaults \emph{of that level}, which
may differ from the defaults of the method. For example, a plain \code{hyperparameters<double>} object passed to
\code{NelderMead_Optimizer} replaces its \code{my_maxstep=5000}, \code{my_f_thr=1e-8}, \code{my_x_thr=1e-6} by 200, 0, 0 (no
convergence test); \code{gradient_hyperparameters} passed to Adam replaces \code{my_maxstep=1000} by 200;
\code{SGD_hyperparameters} passed to QNG brings \code{my_momentum=0.9}. Set those fields explicitly, or use the method's own
class.}

A derived hyperparameter class written for a new method must override \code{setfrom()} in the same way (one line, see the
header), and \code{converged()}, \code{setreached()} if it adds parameters relevant for them.

\subsection{The achieved values}\label{sec:optimize:achieved}

If \code{achieved!=NULL}, at the end \code{achieved->setfrom(hp)} copies the parameters used (at the common level of the two
classes, as above) and \code{achieved->setreached()} overwrites the limits and thresholds with the values reached:
\begin{itemize}
\item \code{my_maxstep}: number of iterations done; \code{my_maxfunc}: number of function evaluations done;
\item \code{my_f_thr}, \code{my_x_thr}: $|\Delta f|$ and $|\Delta x|$ of the last step (0 if no step was done); for
Nelder-Mead the spread of $f$ and the diameter of the simplex, for COBYLA the change of $f$ and the radius $\rho$;
\item \code{my_g_thr} (if \code{achieved} is at least a \code{gradient_hyperparameters}): the last gradient norm
($-1$ if not available);
\item \code{my_saddle_order} (trust-region methods only): the number of negative eigenvalues of the Hessian at the result (of the
exact one if the problem provides \code{OptHessian}, otherwise of the approximation);
\item \code{my_lr} (first-order methods): the initial learning rate actually used (after the SPSA calibration), if
\code{achieved} is at least of the class \code{firstorder_hyperparameters};
\item the other fields: the parameters as they were used.
\end{itemize}
The type of \code{achieved} determines what is returned: an object of the class \code{hyperparameters<double>} suffices for
the counts, the class \code{gradient_hyperparameters<double>} is needed to obtain $|g|$.

%-----------------------------------------------------------------------------------------------------------------------------
\section{Infrastructure and common drivers}\label{sec:optimize:drivers}

This section is needed only for writing new methods; the classes are public so that methods can be added outside the library.

\subsection{Evaluation and iteration bookkeeping}

\begin{lareference}
\lafunc{template<typename T> class optevaluator}{Wraps the problem: counts evaluations and keeps the best point seen.}
\lafunc{const OptProblem<T> *o; int maxfunc; int nfunc; T fbest; NRVec<T> xbest}{The problem, the evaluation budget ($\le0$
unlimited), the number of evaluations done, the lowest $f$ found and its point.}
\lafunc{optevaluator(const OptProblem<T> *oo, const int maxf=0)}{Constructor.}
\lafunc{T operator()(const NRVec<T> &x, NRVec<T> *grad=NULL)}{\code{o->funcgrad(x,grad)}, counted; if $f$ is the lowest so far,
\code{xbest=x} (a shallow copy sharing the data of \code{x}, see section~\ref{sec:optimize:pitfalls}).}
\lafunc{T operator()(const OptResiduals<T> *res, const NRVec<T> &x, NRVec<T> &r)}{Residuals at \code{x} into \code{r}, returns
$\frac12|r|^2$; counted and recorded as above.}
\lafunc{bool exhausted() const; int remaining() const}{Budget used up; number of evaluations left
(\code{std::numeric_limits<int>::max()} if unlimited).}

\lafunc{template<typename T> class optstate}{Iteration bookkeeping common to all methods: convergence and limit tests, verbose
output, result.}
\lafunc{const char *name; const hyperparameters<T> &hp; const optevaluator<T> &e}{Method name for the output, the parameters
used (their virtual \code{converged()} is called), the evaluator.}
\lafunc{int iter; T f; T gnorm; T df; T dx; bool stepdone; bool converged}{Iterations done; current $f$; current $|g|$ ($<0$ if
not available); $|\Delta f|$, $|\Delta x|$ of the last step; whether \code{df}, \code{dx} are valid; whether converged.}
\lafunc{optstate(const char *nm, const hyperparameters<T> &h, const optevaluator<T> &ee)}{Constructor.}
\lafunc{bool stop(const T ff, const T gn= -1, const int needed=1)}{To be called at the beginning of each iteration with the
current $f$ and $|g|$: prints the iteration line (verbose $\ge1$) and returns \code{true} if converged (\code{hp.converged()}),
if \code{iter>=my_maxstep}, or if fewer than \code{needed} evaluations (those of the next iteration) remain.}
\lafunc{void step(const T dff, const T dxx)}{Records a step: $|\Delta f|$, $|\Delta x|$, \code{stepdone=true}.}
\lafunc{optresult<T> finish(hyperparameters<T> *achieved) const}{Final report (verbose $\ge1$), fills \code{achieved} if not
\code{NULL}, returns the result.}
\end{lareference}

\subsection{One-dimensional searches}

\begin{lareference}
\lafunc{template<typename T> T linesearch_wolfe(optevaluator<T> &e, NRVec<T> &x, T &f, NRVec<T> &g, const NRVec<T> &p, const T alpha1, const linesearch_hyperparameters<T> &lp)}{Line
search along the descent direction \code{p} satisfying the strong Wolfe conditions (bracketing and zoom with safeguarded cubic
interpolation, \cite{optimize:nocedal2006} alg.~3.5, 3.6). On input \code{x}, \code{f}, \code{g} at the starting point,
\code{alpha1} the initial trial step; on success they are replaced by the new point and the step length $\alpha$ is returned.
If the strong Wolfe conditions cannot be met within \code{my_maxlinesearch} evaluations (or the remaining budget), the best point
satisfying sufficient decrease is accepted; 0 is returned on failure (also if \code{p} is not a descent direction), \code{x},
\code{f}, \code{g} then unchanged. The step is limited to $|\alpha p|\le$ \code{my_maxsteplength}.}
\lafunc{template<typename T> T linemin_brent(optevaluator<T> &e, NRVec<T> &x, T &f, const NRVec<T> &u, const T t1, const T tol, const T abstol, const int maxiter, const int verbose=0)}{Derivative-free
minimization of $f(x+tu)$: bracketing starting with the trial step \code{t1} (the opposite direction is tried if uphill),
then Brent's method \cite{optimize:brent1973} with relative and absolute tolerances \code{tol}, \code{abstol} on $t$, at most
\code{maxiter} evaluations (and the remaining budget). On output \code{x}, \code{f} at the best point found; returns $t$, or 0
if no decrease was found (\code{x}, \code{f} unchanged).}
\end{lareference}

\subsection{Line-search methods}

\begin{lareference}
\lafunc{template<typename T> class linesearch_direction}{Search direction of a line-search method, abstract.}
\lafunc{virtual NRVec<T> direction(const NRVec<T> &g) =0}{Search direction at the point with gradient \code{g}.}
\lafunc{virtual bool update(const NRVec<T> &s, const NRVec<T> &y, const NRVec<T> &p, const NRVec<T> &g) =0}{After a step
$s=\alpha p$ with gradient change $y$ and new gradient $g$; returns \code{false} if the update was skipped (only reported).}
\lafunc{virtual void reset() =0}{Forget the accumulated information (restart).}
\lafunc{virtual bool steepest() const =0}{No information accumulated, \code{direction()} gives steepest descent.}
\lafunc{virtual bool indefinite() const}{\code{true}: a non-descent direction is legitimate (SR1); the driver then takes a
steepest-descent step without reset. Default \code{false}.}
\lafunc{virtual bool alpha_from_previous() const}{\code{true}: initial trial step $\alpha_{k-1}\,g_{k-1}\cdot p_{k-1}/(g_k\cdot p_k)$
(\cite{optimize:nocedal2006} eq.~3.60, used by CG); otherwise 1 (the quasi-Newton step), or a unit step in $x$ for steepest
descent. Default \code{false}.}
\lafunc{template<typename T> optresult<T> linesearch_minimize(const char *name, linesearch_direction<T> &dir, const OptProblem<T> *o, NRVec<T> &x, const linesearch_hyperparameters<T> &hp, hyperparameters<T> *achieved)}{Common
driver: in each iteration \code{p=dir.direction(g)}; if $p\cdot g\ge0$, steepest descent (with \code{dir.reset()} unless
\code{indefinite()}); \code{linesearch_wolfe()}; on its failure a restart, or the end if it failed along steepest descent;
then \code{dir.update()} and the step is recorded. One evaluation per line search trial point.}
\end{lareference}

\subsection{First-order methods}

\begin{lareference}
\lafunc{template<typename T> class firstorder_step}{Update rule of a first-order method, abstract.}
\lafunc{virtual NRVec<T> step(const NRVec<T> &g, const int k) =0}{Direction $d$ of the step $x\mathrel{-}=\mathrm{lr}\,d$ in
iteration $k=1,2,\dots$ with gradient \code{g}. The gradient vector passed is freshly allocated in each iteration, so the stepper
may keep it.}
\lafunc{template<typename T> class gradient_estimator}{Supplies $f$ and the gradient to the driver; the default implementation
is the exact evaluation, derived classes estimate them (SPSA).}
\lafunc{virtual T operator()(optevaluator<T> &e, const NRVec<T> &x, NRVec<T> &g, const int k)}{Returns $f$ at \code{x} and the
gradient in \code{g}; \code{k=0,1,...} is the iteration.}
\lafunc{virtual int cost() const}{Number of evaluations of one call (default 1).}
\lafunc{virtual bool exact() const}{Default \code{true}. With estimates (\code{false}) no convergence test is done, since small
gradients or changes can occur by chance, and the iterations run until \code{my_maxstep} or \code{my_maxfunc}.}
\lafunc{virtual T calibrate(optevaluator<T> &e, const NRVec<T> &x, const firstorder_hyperparameters<T> &hp)}{Learning rate to be
used; default \code{hp.my_lr}.}
\lafunc{template<typename T> optresult<T> firstorder_minimize(const char *name, firstorder_step<T> &stepper, const OptProblem<T> *o, NRVec<T> &x, const firstorder_hyperparameters<T> &hp, hyperparameters<T> *achieved, gradient_estimator<T> *est=NULL)}{Common
driver: $x\leftarrow x-\mathrm{lr}_k\,d-\mathrm{lr}_k\,\mathrm{weight\_decay}\,x$ with $d$ from the stepper, \code{est==NULL}
meaning exact gradients. With \code{my_natural} the gradient of the estimator is replaced by the natural gradient. Implements
\code{my_return_best}, and returns the learning rate used in \code{achieved} if it is a \code{firstorder_hyperparameters}.}
\end{lareference}

\subsection{Trust-region methods}\label{sec:optimize:tr}

The trust-region methods work with the quadratic model $f+g\cdot p+\frac12 p\cdot Bp$ within $|p|\le\Delta$
\cite{optimize:nocedal2006,optimize:conn2000}. The matrix $B$ is the exact Hessian (\code{OptHessian}) or an approximation
updated by a secant formula.

\begin{lareference}
\lafunc{enum TR_hessian {TR_AUTO, TR_EXACT, TR_BFGS, TR_SR1, TR_PSB, TR_BOFILL}}{Choice of $B$.
\code{TR_EXACT}: exact Hessian (\code{laerror} if the problem does not provide it). \code{TR_AUTO}: exact if available,
otherwise BFGS for minima and Bofill for saddle points. \code{TR_BFGS}, \code{TR_SR1}, \code{TR_PSB}, \code{TR_BOFILL}: updated
by the BFGS, SR1, Powell-symmetric-Broyden or Bofill \cite{optimize:bofill1994} formula starting from the identity, scaled after
the first step by $y\cdot y/s\cdot y$. Between exact Hessians (\code{my_hessian_recompute>1}) the updates are SR1 for minima (the
exact Hessian may be indefinite) and Bofill for saddle points. BFGS is refused for saddle points.}
\lafunc{template<typename T> bool hessian_update(NRSMat<T> &B, const NRVec<T> &s, const NRVec<T> &y, const TR_hessian type, const T thr)}{Update
of a Hessian approximation $B$ (not of its inverse) by the step $s$ and gradient change $y$, $\xi=y-Bs$:
BFGS $B+yy\Tr/(y\cdot s)-Bs(Bs)\Tr/(s\cdot Bs)$ (skipped if $s\cdot y\le\mathrm{thr}\,|s||y|$ or $s\cdot Bs\le0$);
SR1 $B+\xi\xi\Tr/(\xi\cdot s)$ (skipped if $|\xi\cdot s|\le\mathrm{thr}\,|\xi||s|$);
PSB $B+(\xi s\Tr+s\xi\Tr)/(s\cdot s)-(\xi\cdot s)ss\Tr/(s\cdot s)^2$;
Bofill $\varphi\,$SR1$+(1-\varphi)\,$PSB with $\varphi=(\xi\cdot s)^2/(|\xi|^2|s|^2)$ (skipped if $\xi=0$).
Returns \code{false} if skipped. Only \code{TR_BFGS}, \code{TR_SR1}, \code{TR_PSB}, \code{TR_BOFILL} are valid.}

\lafunc{T my_delta0; T my_deltamax}{(\code{trustregion_hyperparameters}) Initial (1) and maximal (\code{1e3}) trust-region
radius.}
\lafunc{T my_eta}{A minimization step is accepted if (actual reduction)/(predicted reduction) $>\eta$ (\code{1e-4}).}
\lafunc{TR_hessian my_hessian}{See above (\code{TR_AUTO}).}
\lafunc{int my_hessian_recompute}{With an exact Hessian: recompute it every $k$-th accepted step and update in between; 1 =
always exact (default).}
\lafunc{T my_update_thr}{Threshold for skipping the updates (\code{thr} of \code{hessian_update()}, \code{1e-8}).}

\lafunc{template<typename T> class trustregion_step}{Step of a trust-region method, abstract.}
\lafunc{virtual NRVec<T> step(const NRVec<T> &g, const NRSMat<T> &B, const T delta) =0}{Step within the trust region.}
\lafunc{virtual bool saddle() const}{Searching for a saddle point: steps are always taken and the radius is controlled by the
quality of the model only. Default \code{false}.}
\lafunc{template<typename T> optresult<T> trustregion_minimize(const char *name, trustregion_step<T> &stepper, const OptProblem<T> *o, NRVec<T> &x, const trustregion_hyperparameters<T> &hp, hyperparameters<T> *achieved)}{Common
driver. For minimization: radius control by \cite{optimize:nocedal2006} alg.~4.1 (ratio $<0.25$: $\Delta=|p|/4$; ratio $>0.75$
and step on the boundary: $\Delta=\min(2\Delta,\Delta_{\max})$); if the predicted change is below the rounding errors of $f$,
the step is judged by the decrease of $|g|$. For saddle points: every step is taken; ratio outside $[0.25,1.75]$ gives
$\Delta=|p|/4$, ratio in $(0.75,1.25)$ with a boundary step doubles $\Delta$. One evaluation (with gradient) per iteration;
rejected steps count as iterations. At the end the Hessian (exact if available, otherwise $B$) is diagonalized and the number of
its negative eigenvalues is printed and returned in \code{achieved->my_saddle_order}.}
\end{lareference}

\subsection{Writing a new method}

A new line-search method needs a \code{linesearch_direction} and an \code{Optimizer} whose \code{minimize()} sets up the
hyperparameters and calls the driver; the same pattern applies to \code{firstorder_step} with the driver
\code{firstorder_minimize()} and to \code{trustregion_step} with the driver \code{trustregion_minimize()}. A complete (compile-tested) example, steepest descent with the
strong Wolfe line search:
\begin{lstlisting}
template<typename T>
class SD_direction : public linesearch_direction<T> {
public:
NRVec<T> direction(const NRVec<T> &g) {NRVec<T> p(g); p*= (T)-1; return p;}
bool update(const NRVec<T> &s, const NRVec<T> &y, const NRVec<T> &p, const NRVec<T> &g) {return true;}
void reset() {}
bool steepest() const {return true;}
};

template<typename T>
class SD_Optimizer : public Optimizer<T> {
public:
optresult<T> minimize(const OptProblem<T> *o, NRVec<T> &x, const hyperparameters<T> *par=NULL, hyperparameters<T> *achieved=NULL)
	{
	linesearch_hyperparameters<T> hp; //a derived class would set the method's defaults
	if(par) hp.setfrom(*par);
	if(hp.my_saddle_order) laerror("SD cannot search for saddle points");
	SD_direction<T> dir;
	return linesearch_minimize("SD",dir,o,x,hp,achieved);
	}
};
\end{lstlisting}
Methods with their own iteration loop use \code{optevaluator} for all evaluations and \code{optstate}: \code{stop()} at the
beginning of each iteration, \code{step()} after each step, \code{finish(achieved)} at the end.

%-----------------------------------------------------------------------------------------------------------------------------
\section{Methods}\label{sec:optimize:methods}

Table~\ref{tab:optimize:methods} summarizes the methods. ``Gradient'' means that \code{funcgrad()} is called with a gradient
request; such methods can be applied to an $f$-only problem through \code{FiniteDifference} or \code{ParameterShift}.
``Optional'' capabilities are used if present.

\begin{table}[htbp]
\centering\footnotesize
\caption{Methods, their requirements and properties.}\label{tab:optimize:methods}
\begin{tabular}{@{}>{\raggedright}p{0.25\textwidth}>{\raggedright}p{0.11\textwidth}>{\raggedright}p{0.24\textwidth}>{\raggedright}p{0.14\textwidth}>{\raggedright\arraybackslash}p{0.12\textwidth}@{}}
\toprule
Class & Deriva\-tives & Capabilities & Constraints, saddles & Noisy $f$ \\
\midrule
\texttt{BFGS\_Optimizer} (BFGS, DFP, SR1) & gradient & -- & no & no \\
\texttt{LBFGS\_Optimizer} & gradient & -- & no & no \\
\texttt{CG\_Optimizer} & gradient & -- & no & no \\
\texttt{GDIIS\_Optimizer} & gradient & -- & no & no \\
\texttt{TrustRegion\_\-Optimizer} & gradient & \texttt{OptHessian} optional (required for \texttt{TR\_EXACT}) & no & no \\
\texttt{PRFO\_Optimizer} & gradient & \texttt{OptHessian} optional & saddles of any order & no \\
\texttt{LM\_Optimizer} & Jacobian optional & \texttt{OptResiduals} required & no & no \\
\texttt{NelderMead\_\-Optimizer} & none & -- & no & partly \\
\texttt{Powell\_Optimizer} & none & -- & no & no \\
\texttt{Rotosolve\_\-Optimizer} & none & \texttt{OptTrigono\-metric} optional & no & with \texttt{my\_re\-evaluate} \\
\texttt{Cobyla\_Optimizer} & none & \texttt{OptConstraints} optional & inequality constraints & no \\
\texttt{SGD\_Optimizer} & gradient & \texttt{OptMetric} if \texttt{my\_natural} & no & yes \\
\texttt{Adam\_Optimizer} & gradient & \texttt{OptMetric} if \texttt{my\_natural} & no & yes \\
\texttt{QNG\_Optimizer} & gradient & \texttt{OptMetric} required & no & yes \\
\texttt{SPSA\_Optimizer} & none & \texttt{OptMetric} if \texttt{my\_natural} & no & yes \\
\bottomrule
\end{tabular}
\end{table}

Every \code{XXX_Optimizer} class below has only the member \code{minimize()} with the signature of
\code{Optimizer<T>::minimize()} and a virtual destructor; only the hyperparameters are listed. All methods except P-RFO call \code{laerror} if
\code{my_saddle_order!=0}.

\subsection{Quasi-Newton methods: BFGS, DFP, SR1}

\code{BFGS_Optimizer} keeps an explicit approximation $H$ of the inverse Hessian, starting from the identity scaled after the
first step by $s\cdot y/y\cdot y$ (\cite{optimize:nocedal2006} eq.~6.20), and uses the strong Wolfe line search with the initial
trial step 1 (the quasi-Newton step). Memory and work are $O(n^2)$ per iteration; for large $n$ use L-BFGS.

\begin{lareference}
\lafunc{template<typename T> class quasinewton_hyperparameters : public linesearch_hyperparameters<T>}{}
\lafunc{T my_sy_thr}{Skip the update if $s\cdot y\le\mathrm{sy\_thr}\,|s||y|$ (curvature condition violated numerically;
\code{1e-12}). Also used by L-BFGS.}
\lafunc{enum QN_update {QN_BFGS, QN_DFP, QN_SR1}}{Update formula of $H$ (\cite{optimize:nocedal2006} eqs.~6.17, 6.15, 6.25).}
\lafunc{template<typename T> class BFGS_hyperparameters : public quasinewton_hyperparameters<T>}{}
\lafunc{QN_update my_update}{\code{QN_BFGS} (default), \code{QN_DFP} (needs an accurate line search, \code{my_c2}$\approx0.1$,
otherwise it can stall), \code{QN_SR1} ($H$ may become indefinite; a non-descent direction then leads to a steepest descent step
without reset). The verbose name is BFGS, DFP or SR1.}
\lafunc{T my_sr1_thr}{SR1: skip the update if $|(s-Hy)\cdot y|\le\mathrm{sr1\_thr}\,|s-Hy||y|$ (\cite{optimize:nocedal2006}
eq.~6.26; \code{1e-8}).}
\lafunc{template<typename T> class BFGS_Optimizer : public Optimizer<T>}{}
\end{lareference}

\subsection{L-BFGS}

Limited-memory BFGS (\cite{optimize:nocedal2006} sec.~7.2): the last \code{my_memory} pairs $(s,y)$ in a ring buffer, the
direction by the two-loop recursion (alg.~7.4) with $H_0=\gamma I$, $\gamma=s\cdot y/y\cdot y$ of the newest pair. Memory
$O(mn)$. Pairs violating the curvature condition (\code{my_sy_thr}) are skipped.
\begin{lareference}
\lafunc{template<typename T> class LBFGS_hyperparameters : public quasinewton_hyperparameters<T>}{}
\lafunc{int my_memory}{Number of stored pairs (10; \code{laerror} if $<1$).}
\lafunc{template<typename T> class LBFGS_Optimizer : public Optimizer<T>}{Verbose name L-BFGS.}
\end{lareference}

\subsection{Nonlinear conjugate gradients}

$p_{k+1}=-g_{k+1}+\beta p_k$ (\cite{optimize:nocedal2006} sec.~5.2, \cite{optimize:hager2005}), with the initial trial step of
the line search from the previous step. Restarts with steepest descent every \code{my_restart} iterations, by Powell's test,
after a failed line search, and if the curvature $p\cdot y\le0$ for the formulas dividing by it.
\begin{lareference}
\lafunc{enum CG_beta {CG_FR, CG_PR, CG_HS, CG_DY, CG_HZ}}{Formula for $\beta$ ($y=g_{k+1}-g_k$):
Fletcher-Reeves $|g_{k+1}|^2/|g_k|^2$; Polak-Ribi\`ere+ $\max(0,g_{k+1}\cdot y/|g_k|^2)$;
Hestenes-Stiefel+ $\max(0,g_{k+1}\cdot y/p_k\cdot y)$; Dai-Yuan $|g_{k+1}|^2/p_k\cdot y$; Hager-Zhang
$\max\bigl(\beta^{\rm HZ},-1/(|p_k|\min(0.01,|g_k|))\bigr)$, $\beta^{\rm HZ}=(g_{k+1}\cdot y-2(p_k\cdot g_{k+1})|y|^2/p_k\cdot y)/p_k\cdot y$.
Verbose names CG-FR, CG-PR+, CG-HS+, CG-DY, CG-HZ.}
\lafunc{template<typename T> class CG_hyperparameters : public linesearch_hyperparameters<T>}{Default \code{my_c2=0.1}
(CG needs a more accurate line search).}
\lafunc{CG_beta my_beta}{\code{CG_PR} (default), see above.}
\lafunc{int my_restart}{Restart every \code{my_restart} iterations, $\le0$ = the dimension $n$ (default 0).}
\lafunc{T my_powell_thr}{Restart if $|g_{k+1}\cdot g_k|\ge\mathrm{powell\_thr}\,|g_{k+1}|^2$ (\cite{optimize:nocedal2006}
eq.~5.52), $\le0$ disables (0.2).}
\lafunc{template<typename T> class CG_Optimizer : public Optimizer<T>}{}
\end{lareference}

\subsection{GDIIS}

Geometry DIIS \cite{optimize:csaszar1984} with a quasi-Newton Hessian \cite{optimize:farkas2002}. The error vectors are the
quasi-Newton steps $e_i=-Hg_i$; the new point is $\sum_ic_i(x_i+e_i)$ with $c$ minimizing $|\sum_ic_ie_i|$ subject to
$\sum_ic_i=1$ (the \code{DIIS} class of diis.h, chapter~\ref{ch:iterative}). One evaluation per iteration. If the DIIS step goes
uphill ($s\cdot g\ge0$), is longer than \code{my_maxsteplength}, or does not decrease $f$, the DIIS history is reset and a
strong Wolfe line search along the quasi-Newton step is done; the first step after a (re)start is always a line search. $H$ is
updated as in \code{BFGS_Optimizer} (all its parameters apply, including \code{my_update}).
\begin{lareference}
\lafunc{template<typename T> class GDIIS_hyperparameters : public BFGS_hyperparameters<T>}{}
\lafunc{int my_diis_dim}{Number of stored points (5; \code{laerror} if $<1$).}
\lafunc{T my_diis_eps}{Regularization added to the diagonal of the DIIS matrix (0).}
\lafunc{template<typename T> class GDIIS_Optimizer : public Optimizer<T>}{}
\end{lareference}

\subsection{Trust-region Newton and quasi-Newton}

Minimization by \code{trustregion_minimize()} (section~\ref{sec:optimize:tr}) with the subproblem
$\min g\cdot p+\frac12p\cdot Bp$, $|p|\le\Delta$, solved either exactly in the eigenbasis of $B$ \cite{optimize:more1983}
(Newton step if $B$ is positive definite and the step lies inside, otherwise $p(\mu)=-(B+\mu I)^{-1}g$ with $|p(\mu)|=\Delta$
found by a safeguarded Newton iteration on $1/|p(\mu)|$, including the hard case; $O(n^3)$ per iteration, handles indefinite
$B$), or approximately by truncated conjugate gradients \cite{optimize:steihaug1983} (\cite{optimize:nocedal2006} alg.~7.2),
which avoids the diagonalization for large $n$. With the exact Hessian (default if \code{OptHessian} is provided) this is a
globally convergent Newton method.
\begin{lareference}
\lafunc{enum TR_subproblem {TR_EIGEN, TR_STEIHAUG}}{Subproblem solver.}
\lafunc{template<typename T> class TrustRegion_hyperparameters : public trustregion_hyperparameters<T>}{All parameters of
\code{trustregion_hyperparameters} (section~\ref{sec:optimize:tr}) and:}
\lafunc{TR_subproblem my_subproblem}{\code{TR_EIGEN} (default) or \code{TR_STEIHAUG}. Verbose names TR-Newton, TR-Steihaug.}
\lafunc{int my_maxcg}{Maximum number of CG iterations of Steihaug's method, $\le0$ = the dimension (0).}
\lafunc{template<typename T> class TrustRegion_Optimizer : public Optimizer<T>}{Minimization only (\code{laerror} for
\code{my_saddle_order!=0}, use P-RFO). \code{achieved->my_saddle_order} returns the Hessian index at the result, which is 0
for a true minimum.}
\end{lareference}
Note that the Steihaug variant works with the stored matrix $B$ (exact or updated), so it saves the $O(n^3)$ diagonalization but
not the $O(n^2)$ storage.

\subsection{P-RFO: minima and saddle points}

Partitioned rational function optimization \cite{optimize:banerjee1985,optimize:baker1986}: search for a stationary point with
\code{my_saddle_order}$=k$ negative Hessian eigenvalues. In the eigenbasis of $B$ ($\lambda_i$ ascending, $\tilde g_i$ the
components of $g$), the step is $\tilde p_i=-\tilde g_i/(\lambda_i-\mu)$ with $\mu=\mu_{\max}$, the largest eigenvalue of the
bordered $(k+1)\times(k+1)$ matrix $\bigl[\begin{smallmatrix}\Lambda_k&\tilde g_k\\ \tilde g_k\Tr&0\end{smallmatrix}\bigr]$ of
the $k$ lowest modes (maximized), and $\mu=\mu_{\min}$, the lowest eigenvalue of the corresponding matrix of the other modes
(minimized); the step is scaled down to $\Delta$ if longer. $k=0$ gives RFO minimization. For $k>0$ the trust-region driver works
in its saddle mode (steps always taken). The modes to be maximized are always the $k$ lowest ones (no mode following).
The Hessian choice follows \code{my_hessian}: with \code{TR_AUTO} the exact Hessian if available, otherwise Bofill updates for
$k>0$ (\code{TR_BFGS} is refused for $k>0$). Convergence is normally by \code{my_g_thr}; the number of negative eigenvalues at
the result is only reported (in the verbose output and in \code{achieved->my_saddle_order}), not checked, so the caller should
compare it with $k$.
\begin{lareference}
\lafunc{template<typename T> class PRFO_hyperparameters : public trustregion_hyperparameters<T>}{All parameters of
\code{trustregion_hyperparameters}, with the default \code{my_delta0=0.3}; \code{my_saddle_order} in $[0,n]$.}
\lafunc{template<typename T> class PRFO_Optimizer : public Optimizer<T>}{Verbose name P-RFO; \code{laerror} for an invalid
saddle order.}
\end{lareference}

\subsection{Levenberg-Marquardt}

For least-squares problems providing \code{OptResiduals} (typically derived from \code{OptLeastSquares}): the step $h$ solves
$(J\Tr J+\mu D)h=-J\Tr r$ (\code{linear_solve}), $D=I$ or $D=\mathrm{diag}(J\Tr J)$ (Marquardt scaling); the damping starts at
$\mu=\tau\max_i(J\Tr J)_{ii}$ (with scaling $\mu=\tau$, the same relative to $D$) and is updated by the gain ratio $\varrho$ \cite{optimize:madsen2004} (alg.~3.16): for $\varrho>0$
the step is accepted and $\mu\leftarrow\mu\max(1/3,1-(2\varrho-1)^3)$, $\nu=2$, otherwise $\mu\leftarrow\mu\nu$, $\nu\leftarrow2\nu$.
The Jacobian comes from \code{jacobian()} of the problem or from forward differences (relative step \code{my_fd_step}, $n$
residual evaluations counted in \code{my_maxfunc}). Convergence: $|J\Tr r|<$ \code{my_g_thr}, the general $f$/$x$ criterion, or
when a rejected step becomes negligible at machine precision ($h_i^2\le(10\epsilon)^2(x_i^2+\epsilon^2)$ for all $i$,
tested per component to be scale invariant, cf.\ \cite{optimize:madsen2004} eq.~3.16). The function value returned is
$f=\frac12|r|^2$. The problem's \code{funcgrad()} is not used.
\begin{lareference}
\lafunc{template<typename T> class LM_hyperparameters : public gradient_hyperparameters<T>}{}
\lafunc{T my_tau}{Initial damping factor $\tau$ (\code{1e-3}).}
\lafunc{bool my_scaling}{Marquardt scaling $D=\mathrm{diag}(J\Tr J)$ (default \code{false}, $D=I$). Recommended for
badly scaled parameters, in float usually necessary (otherwise the steps in the weakly curved directions fall below the
resolution of $f$).}
\lafunc{T my_fd_step}{Relative step of the forward-difference Jacobian (\code{1e-7}).}
\lafunc{template<typename T> class LM_Optimizer : public Optimizer<T>}{\code{laerror} if the problem does not provide
\code{OptResiduals}. Verbose name Levenberg-Marquardt.}
\end{lareference}
\subsection{Nelder-Mead downhill simplex}

Derivative-free simplex method \cite{optimize:lagarias1998} with reflection, expansion, outside and inside contraction and
shrink, starting from the simplex $x_0$, $x_0+\mathrm{initstep}\,e_i$ ($n+1$ evaluations). Convergence: \code{my_f_thr} is
compared with the spread $f_{\rm worst}-f_{\rm best}$ of the function values in the simplex and \code{my_x_thr} with its
diameter (largest distance from the best vertex); both have to be satisfied. The best vertex is returned.
\begin{lareference}
\lafunc{template<typename T> class NelderMead_hyperparameters : public hyperparameters<T>}{Defaults \code{my_maxstep=5000},
\code{my_f_thr=1e-8}, \code{my_x_thr=1e-6}.}
\lafunc{T my_initstep}{Edge length of the initial simplex (0.1).}
\lafunc{bool my_adaptive}{Dimension-dependent coefficients $\alpha=1$, $\gamma=1+2/n$, $\rho=3/4-1/(2n)$, $\sigma=1-1/n$
\cite{optimize:gao2012} (default \code{true}, recommended above about 10 dimensions); otherwise the following ones.}
\lafunc{T my_alpha; T my_gamma; T my_rho; T my_sigma}{Reflection (1), expansion (2), contraction (0.5), shrink (0.5)
coefficients.}
\lafunc{template<typename T> class NelderMead_Optimizer : public Optimizer<T>}{}
\end{lareference}

\subsection{Powell's conjugate direction method}

Derivative-free \cite{optimize:powell1964}: in each iteration (sweep) line minimizations by \code{linemin_brent()} along $n$
directions, initially the coordinate axes; the overall displacement of the sweep replaces the direction of the largest decrease if
Powell's test indicates it is worthwhile, followed by a line minimization along it. Convergence: \code{my_f_thr} and
\code{my_x_thr} are compared with the change of $f$ and $x$ in one sweep, both have to be satisfied. The trial step of each
direction is adapted to its last line minimum.
\begin{lareference}
\lafunc{template<typename T> class Powell_hyperparameters : public hyperparameters<T>}{Defaults \code{my_maxstep=1000},
\code{my_f_thr=1e-10}, \code{my_x_thr=1e-6}.}
\lafunc{T my_initstep}{Initial trial step of the line minimizations (0.1).}
\lafunc{T my_linmin_tol}{Relative tolerance of the line minimizations (\code{1e-4}); the absolute one is $0.01\,$\code{my_x_thr}.}
\lafunc{int my_maxlinemin}{Maximum number of evaluations in one line minimization (50).}
\lafunc{template<typename T> class Powell_Optimizer : public Optimizer<T>}{}
\end{lareference}

\subsection{COBYLA}

Constrained optimization by linear approximations \cite{optimize:powell1994,optimize:powell1998}, derivative-free, minimizing $f$
subject to the inequality constraints of \code{OptConstraints} (without them unconstrained). Linear models of $f$ and of the
constraints are interpolated on a simplex of $n+1$ points (initially $x$, $x+\mathrm{rhobeg}\,e_j$). The trust-region step
minimizes the linear model of $f$ subject to the linearized constraints within $|d|\le\Delta$ (an active-set method on the small
linear program; if the linearized constraints cannot be satisfied within the trust region, the largest violation is minimized
first); the merit function is $f+\mu\cdot$(largest constraint violation) with an automatically increased penalty $\mu$. Geometry
steps keep the simplex acceptable. The resolution $\rho$ decreases from \code{my_rhobeg} to \code{my_rhoend}; the iterations
finish when $\rho$ has reached \code{my_rhoend} and the trust-region step does not improve the point (\code{converged=true}), or
by the general criterion, in which \code{my_f_thr} is compared with the change of $f$ at the best vertex in an iteration and
\code{my_x_thr} with $\rho$ (both 0 by default, i.e.\ off). The best vertex with respect to the merit function is returned;
with verbose output the largest remaining constraint violation is printed. COBYLA does not need a constraint to be satisfied at
the starting point.

\begin{lareference}
\lafunc{template<typename T> class Cobyla_hyperparameters : public hyperparameters<T>}{Default \code{my_maxstep=100000}.}
\lafunc{T my_rhobeg}{Initial trust-region radius, about one tenth of the expected change of $x$ (0.5).}
\lafunc{T my_rhoend}{Final trust-region radius, the accuracy of $x$ (\code{1e-6}).}
\lafunc{bool my_flexible_radius}{\code{true} (default): the trust-region radius $\Delta\ge\rho$ is revised by the reduction
ratio and $\rho$ is reduced by the rules of NEWUOA \cite{optimize:powell2004} (eqs.~7.1, 7.6, 6.16), so that $\Delta$ can also
increase again; geometry steps have the length $\max(\min(0.1\,\mathrm{dist},0.5\Delta),\rho)$. \code{false}: Powell's original
COBYLA with $\Delta=\rho$, $\rho$ halved when reduced.}
\lafunc{template<typename T> class Cobyla_Optimizer : public Optimizer<T>}{}
\end{lareference}

\subsection{Rotosolve}

Sweeps of exact minimizations along single parameters in which $f$ is trigonometric (\code{OptTrigonometric})
\cite{optimize:ostaszewski2021,optimize:nakanishi2020}, for degree $>1$ \cite{optimize:wierichs2022}. For $x_i$ of degree $R$,
$f$ along $x_i$ is reconstructed by a discrete Fourier transform from $2R+1$ equidistant samples over its period (the current
value and $2R$ new evaluations) and minimized without further evaluations (in closed form for $R=1$, otherwise by a dense grid
refined by Brent's method on the reconstruction); the step taken is the shortest one modulo the period. Parameters of degree 0,
or all parameters if the problem does not provide \code{OptTrigonometric}, are minimized by \code{linemin_brent()} as in Powell's
method, whose parameters apply. Convergence as for Powell (change of $f$ and $x$ in one sweep). If the evaluation budget does
not suffice for the next minimization, the sweep is interrupted and the method stops without convergence.
\begin{lareference}
\lafunc{template<typename T> class Rotosolve_hyperparameters : public Powell_hyperparameters<T>}{}
\lafunc{bool my_reevaluate}{Evaluate $f$ at the new point instead of taking the value of the reconstruction (for noisy $f$ or
doubtful degrees; one more evaluation per parameter; default \code{false}).}
\lafunc{template<typename T> class Rotosolve_Optimizer : public Optimizer<T>}{}
\end{lareference}

\subsection{SGD with momentum}

By \code{firstorder_minimize()}: velocity $v\leftarrow\mathrm{momentum}\cdot v+g$ ($v=g$ in the first iteration), step direction
$d=v$, or $d=g+\mathrm{momentum}\cdot v$ (Nesterov) \cite{optimize:sutskever2013}, as in the PyTorch update formula.
\begin{lareference}
\lafunc{template<typename T> class SGD_hyperparameters : public firstorder_hyperparameters<T>}{Default \code{my_lr=1e-2}.}
\lafunc{T my_momentum; bool my_nesterov}{Momentum (0.9; 0 gives plain gradient descent); Nesterov variant (\code{false}).
Verbose names SGD, SGD-Nesterov.}
\lafunc{template<typename T> class SGD_Optimizer : public Optimizer<T>}{}
\end{lareference}

\subsection{Quantum natural gradient}

$x\leftarrow x-\mathrm{lr}\,(G+\mathrm{metric\_reg}\,I)^{-1}g$ \cite{optimize:stokes2020}: SGD without momentum with the natural
gradient; the problem has to provide the metric (\code{OptMetric}), the gradient may come e.g.\ from \code{ParameterShift}.
Momentum, Adam or SPSA can be combined with the natural gradient by setting \code{my_natural=true} in their hyperparameters.
Natural steps are longer than gradient steps with the same learning rate (the metric eigenvalues are typically $\le1/4$): use
lr up to about 0.1 or a larger \code{my_metric_reg}.
\begin{lareference}
\lafunc{template<typename T> class QNG_hyperparameters : public SGD_hyperparameters<T>}{Defaults \code{my_momentum=0},
\code{my_natural=true}, \code{my_lr=0.05}.}
\lafunc{template<typename T> class QNG_Optimizer : public Optimizer<T>}{Forces \code{my_natural=true} also if other
hyperparameters were passed; \code{laerror} if the problem does not provide \code{OptMetric}.}
\end{lareference}

\subsection{Adam, AdamW, AMSGrad}

Adam \cite{optimize:kingma2015} with bias-corrected moment estimates, $d=\hat m/(\sqrt{\hat v}+\epsilon)$; with
\code{my_weight_decay>0} the decoupled weight decay of AdamW \cite{optimize:loshchilov2019}; optionally AMSGrad
\cite{optimize:reddi2018} (the elementwise maximum of the past second moments). \code{my_beta1=0} gives RMSprop.
\begin{lareference}
\lafunc{template<typename T> class Adam_hyperparameters : public firstorder_hyperparameters<T>}{Default learning rate
\code{1e-3} of the base class.}
\lafunc{T my_beta1; T my_beta2}{Decay rates of the first (0.9) and second (0.999) moment estimates.}
\lafunc{T my_eps}{Regularization of the division by the second moment (\code{1e-8}).}
\lafunc{bool my_amsgrad}{AMSGrad (\code{false}). Verbose names Adam, AdamW, AMSGrad.}
\lafunc{template<typename T> class Adam_Optimizer : public Optimizer<T>}{}
\end{lareference}

\subsection{SPSA}

Simultaneous perturbation stochastic approximation \cite{optimize:spall1992,optimize:spall1998}: derivative-free gradient estimate
$\hat g_i=(f(x+c_k\Delta)-f(x-c_k\Delta))/(2c_k\Delta_i)$ with random $\Delta_i=\pm1$ (from \code{randdouble()} of \header{la_random.h}),
two evaluations per estimate independently of the dimension, robust to noise (e.g.\ shot noise in VQE). The function value
reported is the average of the two evaluations. Gain sequences: $a_k=\mathrm{lr}/(1+\mathrm{lr\_decay}\cdot k)^{\mathrm{lr\_power}}$,
i.e.\ Spall's $a/(k+1+A)^\alpha$ with $\mathrm{lr}=a/(1+A)^\alpha$, $\mathrm{lr\_decay}=1/(1+A)$, $\mathrm{lr\_power}=\alpha$;
$c_k=c/(k+1)^\gamma$. Since $f$ and $g$ are estimates, no convergence test is done: the iterations run until \code{my_maxstep}
or \code{my_maxfunc}, and \code{converged} is \code{false}. The first estimate at the starting point is always done, even if
\code{my_maxfunc} is smaller than its $2\cdot$\code{my_resamplings} evaluations. Before it, the learning rate can be calibrated
(evaluations counted; the value used is returned in \code{achieved->my_lr} if \code{achieved} is a
\code{firstorder_hyperparameters}).
\begin{lareference}
\lafunc{template<typename T> class SPSA_hyperparameters : public firstorder_hyperparameters<T>}{Defaults
\code{my_lr_decay=0.01}, \code{my_lr_power=0.602}, \code{my_g_thr=0}.}
\lafunc{T my_c; T my_gamma}{Perturbation size $c$, of the order of the noise of $f$ (0.1), and its decay exponent $\gamma$ (0.101).}
\lafunc{int my_resamplings}{Number of gradient estimates averaged in one iteration (1; \code{laerror} if $<1$).}
\lafunc{int my_calibrate}{Number of gradient estimates at the starting point for the calibration of lr, 0 = no calibration (20).}
\lafunc{T my_target_step}{Calibration: lr is set so that the first step has on average this length (0.1).}
\lafunc{template<typename T> class SPSA_Optimizer : public Optimizer<T>}{}
\end{lareference}

%-----------------------------------------------------------------------------------------------------------------------------
\section{Examples}\label{sec:optimize:examples}

All examples were compiled and run; they assume \texttt{\#include "la.h"} and \code{using namespace LA;}.

\subsection{A user-defined problem and BFGS}

The extended Rosenbrock function with its analytic gradient. Note the \code{copyonwrite()} on the gradient vector before its
elements are written.
\begin{lstlisting}
class Rosenbrock : public OptProblem<double> {
public:
double funcgrad(const NRVec<double> &x, NRVec<double> *g) const
	{
	const int n=x.size();
	double f=0;
	if(g) {g->copyonwrite(); *g=0.;} //g may share data with another vector
	for(int i=0; i<n-1; ++i)
		{
		double a=1-x[i], b=x[i+1]-x[i]*x[i];
		f+=a*a+100*b*b;
		if(g) {(*g)[i]+= -2*a-400*x[i]*b; (*g)[i+1]+=200*b;}
		}
	return f;
	}
};

Rosenbrock problem;
NRVec<double> x(-1.2,4); //starting point (-1.2,...,-1.2)
BFGS_hyperparameters<double> param;
param.my_g_thr=1e-8;
param.my_maxstep=1000;
param.my_verbose=0;
gradient_hyperparameters<double> achieved;
BFGS_Optimizer<double> o;
optresult<double> res=o.minimize(&problem,x,&param,&achieved);
std::cout <<"converged "<<res.converged<<" f= "<<res.f<<" x= "<<x
	<<" iterations "<<achieved.my_maxstep<<" evaluations "<<achieved.my_maxfunc
	<<" |g| "<<achieved.my_g_thr<<std::endl;
\end{lstlisting}
(Here \code{x} is read in \code{funcgrad()} through a const reference, so no DEBUG count check is triggered.) The same problem
can be given to any gradient method by changing the optimizer and hyperparameter classes, and to derivative-free methods
directly. With \code{FiniteDifference<double> fd(&problem);} an $f$-only problem can be minimized by BFGS as
\code{BFGS_Optimizer<double>().minimize(&fd,x);}.

\subsection{Freezing parameters}\label{sec:optimize:ex:frozen}

\begin{lstlisting}
Rosenbrock problem;
NRVec<double> x(0.5,6);
bitvector frozen(6);
frozen.clear();
frozen.set(0); frozen.set(3); //x_0 and x_3 keep their values
FrozenOptProblem<double> fp(&problem,x,frozen);
NRVec<double> xr=fp.reduce(x); //the 4 free parameters
LBFGS_Optimizer<double> o;
double f=o.minimize(&fp,xr); //default parameters; optresult converts to f
NRVec<double> xopt=fp.expand(xr); //full vector, x_0 = x_3 = 0.5
\end{lstlisting}
\code{fp.xfull} shares its data with \code{x}; reading elements of \code{x} afterwards in a DEBUG build has to go through a const
reference, writing them needs \code{x.copyonwrite()} first (which then also detaches \code{x} from \code{fp.xfull}).

\subsection{COBYLA with constraints}

\begin{lstlisting}
//min (x0-1)^2+(x1-2)^2 subject to x0+x1 <= 1, solution (0,1)
class ConstrainedQuadratic : public OptProblem<double>, public OptConstraints<double> {
public:
double funcgrad(const NRVec<double> &x, NRVec<double> *g) const
	{return (x[0]-1)*(x[0]-1)+(x[1]-2)*(x[1]-2);} //COBYLA never asks for g
int nconstraints() const {return 1;}
void constraints(const NRVec<double> &x, NRVec<double> &c) const
	{c.copyonwrite(); c[0]=1-x[0]-x[1];} //c_i >= 0
};

ConstrainedQuadratic problem;
NRVec<double> x(0.,2);
Cobyla_hyperparameters<double> param;
param.my_rhobeg=0.5;
param.my_rhoend=1e-8;
param.my_verbose=0;
Cobyla_Optimizer<double> o;
optresult<double> res=o.minimize(&problem,x,&param);
NRVec<double> c(1);
problem.constraints(x,c); //check the feasibility
\end{lstlisting}

\subsection{Least-squares fit by Levenberg-Marquardt}

\begin{lstlisting}
//fit of y = b0 exp(-b1 t) to data; residuals r_i = model - y_i
class ExpFit : public OptLeastSquares<double> {
public:
NRVec<double> t,y; //data
ExpFit(const NRVec<double> &tt, const NRVec<double> &yy) : t(tt), y(yy) {};
int nresiduals() const {return t.size();}
void residuals(const NRVec<double> &b, NRVec<double> &r) const
	{
	r.copyonwrite();
	for(int i=0; i<t.size(); ++i) r[i]=b[0]*std::exp(-b[1]*t[i])-y[i];
	}
bool jacobian(const NRVec<double> &b, NRMat<double> &J) const
	{
	J.copyonwrite();
	for(int i=0; i<t.size(); ++i)
		{
		const double e=std::exp(-b[1]*t[i]);
		J(i,0)=e; J(i,1)= -b[0]*t[i]*e;
		}
	return true; //false would mean: use finite differences
	}
};

const int m=10;
NRVec<double> t(m), y(m);
for(int i=0; i<m; ++i) {t[i]=0.5*i; y[i]=3*std::exp(-0.7*t[i])+0.01*((i%3)-1);}
ExpFit problem(t,y);
NRVec<double> b(2);
b[0]=1; b[1]=0.1;
LM_hyperparameters<double> param;
param.my_g_thr=1e-10;
param.my_verbose=0;
hyperparameters<double> achieved;
LM_Optimizer<double> o;
optresult<double> res=o.minimize(&problem,b,&param,&achieved);
//res.f = 1/2 |r|^2, i.e. half the residual sum of squares
\end{lstlisting}
Since \code{ExpFit} is an \code{OptLeastSquares}, the same object can also be passed to e.g.\ \code{BFGS_Optimizer}.

\subsection{Further examples in the distribution}

t.cc contains testbeds (blocks starting with \code{//testbed of the opt...}) for the methods on \code{testproblem},
\code{testproblem_vqe(_constrained)} (with \code{ParameterShift}), \code{FrozenOptProblem}, P-RFO on the M\"uller-Brown potential
and LM on a NIST problem. The header testoptimize.h (in the source distribution, not installed) defines these test problems, which can serve as
templates for user problems:
\begin{itemize}
\item \code{testproblem<T>}: a smooth function with analytic gradient;
\item \code{testproblem_vqe<T>}: a VQE-like Ising chain in a product state providing \code{OptHessian}, \code{OptMetric} and
\code{OptTrigonometric}; \code{testproblem_vqe_constrained<T>}: the same with one constraint (\code{OptConstraints});
\item \code{testproblem_mullerbrown<T>}: the M\"uller-Brown potential \cite{optimize:mueller1979} with analytic Hessian, three
minima and two saddle points;
\item NIST StRD nonlinear regression problems (\code{OptLeastSquares}) with starting values, certified parameters and
residual sums of squares: \code{testproblem_misra1a<T>}, \code{testproblem_mgh09<T>} (analytic Jacobians) and
\code{testproblem_thurber<T>} (finite-difference Jacobian).
\end{itemize} The regression program testoptimize.cc, run by
\code{make check}, tests all methods, adapters, capabilities, the \code{maxfunc} limits, the hyperparameter inheritance and the
copy-on-write behaviour of the functions used; it prints PASS/FAIL lines and its exit code is the number of failures.

%-----------------------------------------------------------------------------------------------------------------------------
\section{Pitfalls}\label{sec:optimize:pitfalls}

\lawarning{\textbf{Gradient and other output arguments.} The vector \code{*grad} passed to \code{funcgrad()}, and the arguments
\code{h}, \code{g}, \code{c}, \code{r}, \code{J} of the capability functions, are allocated by the caller but may share their
data with other objects. Call \code{copyonwrite()} before writing elements through \code{operator[]}/\code{operator()}
(assignments of whole objects or scalars, e.g.\ \code{*grad=0.;}, are safe and detach by themselves). Forgetting it can silently
corrupt vectors of the optimizer (e.g.\ the previous gradient).}

\lawarning{\textbf{Shared reference to the best point.} \code{optevaluator} keeps \code{xbest=x}, a shallow copy sharing the data
of the vector evaluated. Code calling an \code{optevaluator} (a new method, or a gradient estimator) and then modifying that vector
element-wise must call \code{copyonwrite()} before each such modification, otherwise the recorded best point changes as well
(see \code{forward_difference_jacobian} and \code{trig_minimize} in optimize.cc). Whole-vector operations (\code{axpy},
\code{+=}, assignment) are safe.}

\lawarning{\textbf{Reading possibly shared vectors.} After \code{minimize()} the vector \code{x}, and the vectors given to
\code{FrozenOptProblem} (\code{xfull}), may share data with other objects. In DEBUG builds the non-const \code{operator[]} raises
\code{laerror} when the reference count is $>1$ (section~\ref{sec:concepts:errors}); read elements of a possibly shared non-const
vector through a const reference, e.g.\ \code{const NRVec<double> &xc=x; ... xc[i]}. Inside \code{funcgrad()} the argument
\code{x} is already const.}

\lawarning{\textbf{Convergence thresholds.} The general criterion needs \emph{both} \code{my_f_thr} and \code{my_x_thr} (strict
inequalities); with the default 0 it is off, and gradient methods then converge only by \code{my_g_thr}. When passing a
hyperparameter object of a more general class, its defaults replace those of the method (section~\ref{sec:optimize:setfrom}).
Methods print by default (\code{my_verbose=1}); set it to 0 for silent operation.}

\lawarning{\textbf{Saddle points.} \code{converged} of P-RFO does not imply the requested saddle order; check
\code{achieved->my_saddle_order}. Likewise a converged \code{TrustRegion_Optimizer} run may end at a saddle point if started at
one with zero gradient components along the negative modes; the reported index tells.}

\lawarning{\textbf{Noisy functions.} SPSA (and first-order methods with an inexact \code{gradient_estimator}) never report
convergence; set \code{my_maxstep} or \code{my_maxfunc}. \code{my_return_best} should not be used with noisy $f$. SPSA uses the
global random generator of \header{la_random.h}; seed it for reproducible runs.}
